\documentclass[11pt]{article}
\usepackage[margin=1in]{geometry}
\usepackage{amsmath,amssymb,amsthm,mathtools,bm}
\usepackage{microtype}
\usepackage[T1]{fontenc}
\usepackage{lmodern}
\usepackage{booktabs,enumitem}
\usepackage[hidelinks]{hyperref}
\setlist{nosep}
\newtheorem{theorem}{Theorem}[section]
\newtheorem{corollary}[theorem]{Corollary}
\theoremstyle{definition}
\newtheorem{definition}[theorem]{Definition}
\newtheorem{assumption}[theorem]{Assumption}
\theoremstyle{remark}
\newtheorem{remark}[theorem]{Remark}
\newtheorem{conjecture}{Conjecture}

\newcommand{\Mb}{M_{\mathrm{b}}}
\newcommand{\vf}{v_{\mathrm{f}}}
\newcommand{\qeff}{q_{\mathrm{eff}}}
\newcommand{\FH}{F_H}
\newcommand{\dd}{\mathrm{d}}
\newcommand{\Tr}{\operatorname{Tr}}

\title{\textbf{Synopsis: Hysterical Response Halo and Galactic Scaling}\\[0.35em]
\large What the theorems say (without the proofs)}
\author{Andrew Bruce Boyd}
\date{September 2026}

\begin{document}
\maketitle
\noindent\textit{Computational note.}
Proofs, exploratory experiments, and paper writing were assisted by generative AI.
Mathematical theorems stated in this paper are formally verified in Lean~4 in the
companion specifications, as faithful ports of the TeX arguments at the abstractions
recorded there, except results explicitly tagged as \emph{literature hypotheses}
(standard theorems cited from the literature and not re-proved here) or
\emph{physical inputs} (modeling assumptions outside mathematics).

\begin{abstract}
This is a short guide to the main claims of the galactic modelling paper.
It builds on the Hysterical Response of the earlier response-theory paper.
Each section says \emph{why} the claim matters, then states the result.
The proofs live in the full paper.
\end{abstract}

\section{From Hysterical Response to a galactic halo}
The earlier response-theory paper isolates the Hysterical Response: the residual force of an already-existing interaction under exponentially stable open-system dynamics,
\begin{equation}
  \FH=-\Tr\!\bigl(F\,\mathcal L^{-1}\mathcal S\bigr),
\end{equation}
equivalently $\delta\rho_{\mathrm{ss}}=-\mathcal L^{-1}\mathcal S$.
No new force carrier is introduced: $\FH$ is a component of the response of the ordinary force~$F$ itself.
That paper stops at channel geometry, neutrality, instability, and infrared equivalence.
It does \emph{not} build galaxies.

The present paper takes that gravitational Hysterical Response as the force and continues it into an explicit spatial halo.
A ballistic Wigner representation of the residual response (production source $S$, loss rate $\Gamma$) is phase-space bookkeeping for how the residue is organized in space---not a claim of a new messenger particle.
Here $S$ is the ballistic image of the companion paper's residual drive~$\mathcal S$.
The retarded characteristic solution of $v\cdot\nabla f=S-\Gamma f$ is
$f(x,v)=\int_0^\infty e^{-\Gamma t}S(x-tv,v)\,dt$ (Lean-checked under standard regularity).
The total production rate $Q=\int J$ becomes the amplitude that later sets galactic scales.
If that ballistic $J$ packages Paper~1's unentangled-incoming drive, then $Q=\alpha J_{\mathrm{in}}^{\mathrm{tot}}$; when baryonic accretion supplies those arrivals, $Q=\alpha\dot N_{\mathrm{acc}}$ (conversion $\alpha$ remains model-dependent---not the unit identity $Q=\dot\Mb$).
``Hysterical Response Halo'' means that extended residual-response density (and the acceleration it produces under the infrared closure)---ordinary gravity's leftover, not a new force law invented here from scratch.

\section{Two galactic problems}
Observed rotation curves stay roughly flat where Newtonian baryons would already be Keplerian, and the asymptotic speed tracks baryonic mass as $\Mb\propto\vf^{4}$ (BTFR).
Those are different requirements: one is a radial profile $g_{\mathrm{resp}}\propto1/r$; the other is a mass normalization.
The full paper separates them.

\section{Ballistic response halo}
Under an infrared ballistic realization---local velocity isotropy plus a single characteristic speed~$v$---residual-response production takes the monoenergetic shell form $S=J\,\delta(|u|-v)/(4\pi v^{2})$.
For a localized, centered spatial profile~$J$, the stationary density then has a universal leading far field.

\begin{theorem}[Ballistic response halo]
Outside the source and well inside the propagation length,
\[
  n=\frac{Q}{4\pi v r^{2}}+O(r^{-4}).
\]
The leading halo is spherical; source shape enters only at higher order.
\end{theorem}

\section{Why $g_{\mathrm{resp}}=CQ/r$}
A density $n\propto Q/r^{2}$ is not yet a force; the force is already the Hysterical Response of gravity.
The primary modelling closure is the linear Poisson law for the associated response potential
\begin{equation}
  \nabla^{2}\Phi_{\mathrm{resp}}=\kappa n.
\end{equation}
Together with the ballistic halo and $\nabla^{2}\ln r\propto r^{-2}$, this yields
\begin{equation}
  g_{\mathrm{resp}}=\frac{CQ}{r}
\end{equation}
for a coupling $C=\kappa/(4\pi v)>0$.
Separate-field or unified-field language may be used as brief physical readings of that same infrared law; they do not independently postulate $CQ/r$.

\begin{theorem}[Flat response rotation]
Circular balance $v_{c}^{2}/r=CQ/r$ then implies $v_{c}^{2}=CQ$, hence an asymptotic response-dominated speed $\vf^{2}=CQ$.
\end{theorem}

Flat curves are therefore the radial consequence of the Hysterical Response halo plus the Poisson closure; they do not by themselves fix how $Q$ scales with baryonic mass.

\section{Capacity attractor vs formation conjecture}
The ballistic spatial model fixes the shape of $n(r)$, not the pair $(Q,\Mb)$.
Introduce a capacity coefficient $A>0$ and filling fraction $y=A^{2}\Mb/Q^{2}$
(equivalently the capacity inequality $\Mb\le Q^{2}/A^{2}$ when $y\le1$).
Response-limited assembly is modelled by a logistic ODE for $y$ with stable fixed point $y_{\ast}\in(0,1]$.

\begin{corollary}[Square-root attractor]
\emph{Given} that capacity inequality / filling setup and the stable assembly dynamics, at $y_{\ast}$ one has
\begin{equation}
  Q^{2}=\frac{A^{2}}{y_{\ast}}\Mb,
\end{equation}
so $Q\propto\sqrt{\Mb}$.
This attractor is a theorem/corollary of those hypotheses---not a conjecture.
\end{corollary}

Combined with $\vf^{2}=CQ$,
\begin{equation}
  \vf^{4}=C^{2}Q^{2}\propto\Mb.
\end{equation}

\begin{corollary}[BTFR exponent]
On that attractor, $\vf^{4}\propto\Mb$.
The zero point remains encoded in $C$, $A$, and $y_{\ast}$; those microscopic values are not derived here.
\end{corollary}

\begin{conjecture}[Formation / retention capacity]
What \emph{is} conjectural is the physical selection rule: galactic formation and mature retention require
\[
  \boxed{Q\gtrsim O\bigl(\sqrt{\Mb}\bigr)}.
\]
That requirement is not proved from the ballistic spatial model (the attractor corollary places mature systems \emph{on} $Q\sim\sqrt{\Mb}$ once capacity filling is assumed).
\end{conjecture}

\section{Late-time size--mass law}
With isotropic accretion ($j\propto\Mb^{2/3}$) and size proxy $R\sim j/\vf$, the BTFR velocity scaling implies a definite size--mass exponent.
The isotropic-$j$ input is physically motivated by late-time feeding from a homogeneous intergalactic dust/gas reservoir (no preferred inflow direction); filamentary anisotropy would change the exponent.

\begin{theorem}[Late-time size--mass]
On the mature attractor $\vf^{4}=K\Mb$,
\[
  R\propto\Mb^{5/12}.
\]
\end{theorem}

Observed late-type slopes typically range from $\sim0.22$ (CANDELS late-types) through $\sim0.39$ (high-mass SDSS late-types) to $\sim0.385$ (baryonic disc scale lengths).
The prediction $5/12\approx0.42$ sits near the steeper end of that range.

\section{Asymptotic oblateness}
The first nonspherical correction for a planar axisymmetric source is a transport quadrupole.
Under the same primary Poisson closure it produces an oblate potential whose effective flattening satisfies
\[
  \qeff^{-2}(R)=1+\frac{s}{R^{2}}+O(R^{-4}),
\]
with $s=\langle R'^{2}\rangle$ (and $s=6R_{d}^{2}$ for an exponential disc).
Thus the response is oblate at finite radius and becomes spherical as $R\to\infty$.

\section{What is theorem vs assumption}
\begin{itemize}
\item From the earlier Hysterical Response paper: $\FH=-\Tr(F\mathcal L^{-1}\mathcal S)$ of ordinary gravity; no new force carrier.
\item IR ballistic shell: local velocity isotropy + single IR scale $\Rightarrow$ $S=J\,\delta(|u|-v)/(4\pi v^{2})$.
\item Ballistic spatial model of that residue: $r^{-2}$ halo and quadrupole (continuum asymptotics tagged where Lean does not re-prove the integral expansion).
\item Modelling closure: primary Poisson $\nabla^{2}\Phi_{\mathrm{resp}}=\kappa n$ $\Rightarrow$ $g_{\mathrm{resp}}=CQ/r$ (and the same map for $\qeff$); separate/unified as readings only.
\item Transport PDE: retarded characteristic solution solves $v\cdot\nabla f=S-\Gamma f$ (Lean).
\item ODE / attractor (theorem): given capacity inequality and filling dynamics, stable $y_{\ast}$ and $Q^{2}=(A^{2}/y_{\ast})\Mb$, hence BTFR exponent.
\item Accretion $\to Q$ (response-theory link): $Q=\alpha J_{\mathrm{in}}^{\mathrm{tot}}=\alpha\dot N_{\mathrm{acc}}$ when accretion supplies unentangled arrivals; not the unit identity $Q=\dot\Mb$.
\item Conjecture (not from the spatial model): formation/retention requires $\boxed{Q\gtrsim O(\sqrt{\Mb})}$.
\item Physical inputs (angular momentum): isotropic $j\propto\Mb^{2/3}$ (homogeneous intergalactic dust/gas) and $R\sim j/\vf$ give $a=5/12$.
\item Not claimed: new force carriers; microscopic $C,A,\alpha$; absolute BTFR zero point; Galactic numerical fits; proof of the formation requirement from the spatial model.
\end{itemize}

\end{document}
